\section{Classification and tameness}\label{sec:classification}

We now pass from the dependent pair in Theorem~\ref{thm:dependent-pair} to a canonical form. The classification rests on a polynomial calculation over the original field.

\begin{theorem}[Classification]\label{thm:classification}
Let $H\in K[x,y,z]^3$ satisfy $H(0)=0$, and suppose that its components are linearly independent over $K$. Then $JH$ is nilpotent if and only if there exist $T\in\GL_3(K)$ and a nonconstant polynomial $P\in K[t]$, with $P(0)=0$, such that
\begin{equation}\label{eq:canonical}
 T\circ H\circ T^{-1}=\Ncal_P,
 \qquad
 \Ncal_P=\bigl(P(y+x^2),\ z-2xP(y+x^2),\ P(y+x^2)^2\bigr).
\end{equation}
\end{theorem}

The theorem describes a family of normal forms; it does not assert that $P$ is unique under linear conjugation. We prove the algebraic part separately for use in the descent in Section~\ref{sec:dependent-pair}.

\subsection{From a dependent pair to a quadratic coordinate}
\begin{proposition}\label{prop:pair-normal}
Let $H\in K[x,y,z]^3$ be normalized, with nilpotent Jacobian and linearly independent components. If two independent constant linear combinations of its components are algebraically dependent over $K$, then $H$ is linearly conjugate over $K$ to $\Ncal_P$ for some $P\in K[t]\setminus K$ with $P(0)=0$.
\end{proposition}

\begin{proof}
A constant conjugation and the closed-generator theorem give
\begin{equation}\label{eq:pair-form}
 H=(f(r),v,g(r)),\qquad r(0)=f(0)=g(0)=0,
\end{equation}
where $r\in K[x,y,z]$ is closed and $f,g\in K[t]$ are linearly independent. Their derivatives are likewise linearly independent, so
\begin{equation}\label{eq:a-ratio}
 a(t)=\frac{g'(t)}{f'(t)}\in K(t)\setminus K,\qquad a'(t)\neq0.
\end{equation}
Primes in this proof denote univariate derivatives.

Put $W=f'(r)r_x+g'(r)r_z$. The determinant vanishes because the first and third Jacobian rows are proportional. The other nilpotence equations are
\begin{equation}\label{eq:pair-nilpotence}
 W+v_y=0,\qquad Wv_y-r_y\bigl(f'(r)v_x+g'(r)v_z\bigr)=0.
\end{equation}
We claim that $r_y\neq0$. If $r_y=0$, the two equations give $W=v_y=0$. Reducing $W=0$ modulo $r-\lambda$ gives
\[
 r-\lambda\mid f'(\lambda)r_x+g'(\lambda)r_z.
\]
The polynomial on the right has degree less than $\deg r$, and hence is zero. Two choices of $\lambda\in K$ with independent vectors $(f'(\lambda),g'(\lambda))$ force $r_x=r_z=0$, a contradiction.

\smallskip\noindent\emph{A derivation with a slice.}
On $\mathcal K=K(x,y,z)$ define
\begin{equation}\label{eq:delta}
 \Delta=\partial_x+a(r)\partial_z+\frac{v_y}{f'(r)r_y}\partial_y.
\end{equation}
Equations~\eqref{eq:pair-nilpotence} give
\[
 \Delta x=1,\qquad \Delta r=\Delta v=0,\qquad \Delta z=a(r).
\]
Set $s=z-a(r)x$. Then $\Delta s=0$, and $\det J(x,r,s)=r_y\neq0$. Thus $\mathcal K/K(x,r,s)$ is finite separable. Extend $\partial_r,\partial_s$ to $\mathcal K$. They commute with $\Delta$, the extension of $\partial_x$ in these coordinates. If $M=\ker_{\mathcal K}\Delta$, then both preserve $M$.

Evaluating on $x,r,s$ and using uniqueness of extension gives
\begin{equation}\label{eq:dy-identity}
 \partial_y=r_y\bigl(\partial_r-a'(r)x\partial_s\bigr).
\end{equation}
Since $v\in M$,
\[
 \Delta y=\frac{v_r-a'xv_s}{f'},
\]
where $a',f'$ are evaluated at $r$. Subtracting a polynomial in $x$ and applying $\Delta$ gives
\begin{equation}\label{eq:y-polynomial}
 y=c_0+\frac{v_r}{f'}x-\frac{a'v_s}{2f'}x^2,\qquad c_0\in M.
\end{equation}
Together with $z=s+a(r)x$, this proves $K[x,y,z]\subseteq M[x]$. The variable $x$ is transcendental over $M$: differentiating a minimal polynomial would otherwise contradict $\Delta x=1$ and separability.

\smallskip\noindent\emph{A unit calculation.}
Apply~\eqref{eq:dy-identity} to $y$:
\begin{equation}\label{eq:unit-identity}
 1=r_y\bigl(\partial_r-a'x\partial_s\bigr)y.
\end{equation}
Both factors are in $M[x]$. For the second factor, use~\eqref{eq:y-polynomial} and the fact that the extended derivations preserve $M$. Both factors are therefore units of $M[x]$. The coefficient of $x^3$ in the second factor is
\[
 \frac{(a')^2}{2f'}v_{ss},
\]
and therefore $v_{ss}=0$. The elements $v_s$ and $v-sv_s$ are killed by $\Delta$ and $\partial_s$. Lemma~\ref{lem:field-facts} shows that they are algebraic over $K(r)$. Since $r$ is closed, they lie in $K(r)$. Consequently
\begin{equation}\label{eq:v-rational}
 v=\alpha(r)z+\beta(r)x+\gamma(r),\qquad \alpha,\beta,\gamma\in K(t).
\end{equation}
More precisely, $\beta=-a\alpha$, although this relation is not needed in the remaining reduction.

We next show that the coefficients in~\eqref{eq:v-rational} are polynomials. Write
\[
 d(r)v=\alpha_0(r)z+\beta_0(r)x+\gamma_0(r)
\]
with $d,\alpha_0,\beta_0,\gamma_0\in K[t]$ have greatest common divisor one. If $d$ is nonconstant, choose a root $\lambda$ over $\overline K$. Reduction modulo $r-\lambda$ shows that $r-\lambda$ divides
$\alpha_0(\lambda)z+\beta_0(\lambda)x+\gamma_0(\lambda)$. The divisor has positive $y$-degree, whereas the dividend is independent of $y$. The dividend is therefore zero. All four univariate polynomials then share the root $\lambda$, contradicting their B\'ezout identity. Thus
\begin{equation}\label{eq:v-polynomial}
 \alpha,\beta,\gamma\in K[t].
\end{equation}

\smallskip\noindent\emph{Commuting affine derivations.}
Substitute~\eqref{eq:v-rational} into the trace equation. We obtain
\begin{equation}\label{eq:affine-trace}
 f'(r)r_x+g'(r)r_z+
 \bigl(\alpha'(r)z+\beta'(r)x+\gamma'(r)\bigr)r_y=0.
\end{equation}
For $\lambda\in K$, let
\[
 E_\lambda=f'(\lambda)\partial_x+g'(\lambda)\partial_z+
 \bigl(\alpha'(\lambda)z+\beta'(\lambda)x+\gamma'(\lambda)\bigr)\partial_y.
\]
Reduction of~\eqref{eq:affine-trace} modulo $r-\lambda$ gives $r-\lambda\mid E_\lambda r$. Since $\deg E_\lambda r\leq\deg r$, we have $E_\lambda r=\mu_\lambda(r-\lambda)$ for some $\mu_\lambda\in K$. Assign weights $1,2,1$ to $x,y,z$. The derivation $E_\lambda$ strictly lowers weighted degree, so it is locally nilpotent. This forces $\mu_\lambda=0$, so $E_\lambda r=0$ for every $\lambda\in K$.

Choose two values of $\lambda$ for which $(f'(\lambda),g'(\lambda))$ are independent. Constant linear combinations of the corresponding derivations give
\[
 D_1=\partial_x+b_1(x,z)\partial_y,\qquad
 D_2=\partial_z+b_2(x,z)\partial_y,
\]
with $b_1,b_2$ affine and $D_1r=D_2r=0$. Their commutator is
\[
 [D_1,D_2]=(\partial_xb_2-\partial_zb_1)\partial_y.
\]
Its coefficient is constant. Since the commutator annihilates $r$ and $r_y\neq0$, its coefficient is zero. There is therefore a polynomial $Q\in K[x,z]$ with
\[
 \deg Q\leq2,\quad Q(0)=0,\qquad Q_x=-b_1,\quad Q_z=-b_2.
\]
In the auxiliary polynomial coordinates $(x,\rho,z)$, where $\rho=y+Q(x,z)$, the derivations become $\partial_x,\partial_z$. Their common kernel in the polynomial ring is $K[\rho]$, so $r=R(\rho)$ for a polynomial $R$ with $R(0)=0$. Absorb $R$ into the univariate polynomials in~\eqref{eq:pair-form} and~\eqref{eq:v-rational}.

Write
\[
 Q=\tfrac12Ax^2+Bxz+\tfrac12Cz^2+Dx+Ez.
\]
Expressing the trace equation in the coordinates $x,\rho,z$ and comparing coefficients gives
\[
 \beta'=-Af'-Bg',\qquad \alpha'=-Bf'-Cg',\qquad \gamma'=-Df'-Eg'.
\]
Integrating and using $H(0)=0$, we obtain
\begin{equation}\label{eq:quadratic-form}
 H=\bigl(f(\rho),-Q_xf(\rho)-Q_zg(\rho)+\Lambda(x,z),g(\rho)\bigr),
 \qquad \rho=y+Q(x,z),
\end{equation}
where $\Lambda$ is a homogeneous linear polynomial and $f(0)=g(0)=0$. If $Q_2$ is the quadratic homogeneous part of $Q$, expansion of the principal minors gives
\begin{align}\label{eq:quadratic-relation-derivative}
 \sigma_2(JH)
 &=f'(\rho)\bigl(Q_{xx}f(\rho)+Q_{xz}g(\rho)-\Lambda_x\bigr)\notag\\
 &\quad+g'(\rho)\bigl(Q_{xz}f(\rho)+Q_{zz}g(\rho)-\Lambda_z\bigr)\\
 &=\left.\frac{\dd}{\dd t}
 \bigl(Q_2(f(t),g(t))-\Lambda(f(t),g(t))\bigr)\right|_{t=\rho}.\notag
\end{align}
Substitution $t\mapsto \rho$ is injective. The polynomial being differentiated is constant and vanishes at zero. Thus
\begin{equation}\label{eq:quadratic-relation}
 Q_2(f(t),g(t))=\Lambda(f(t),g(t)).
\end{equation}

\smallskip\noindent\emph{Constant linear normalization.}
The shear $(x,y,z)\mapsto(x,y+Dx+Ez,z)$ removes the linear part of $Q$ from~\eqref{eq:quadratic-form}; the corresponding target change cancels the terms $-Df-Eg$. We may therefore assume $Q=Q_2$. If $Q_2=0$, relation~\eqref{eq:quadratic-relation} and independence of $f,g$ force $\Lambda=0$. The second component would then vanish, contradicting independence.

The binary quadratic form $Q_2$ cannot have rank two. Over $\overline K$, such a form is a product of independent linear forms. After evaluation at $(f,g)$, relation~\eqref{eq:quadratic-relation} would give
\[
  \bar{F}G=\alpha \bar{F}+\beta G,\qquad \bar{F}(0)=G(0)=0,
\]
with $\bar{F},G$ linearly independent over $ \overline K$. But
$(\bar{F}-\beta)(G-\alpha)=\alpha\beta$. If $\alpha\beta\neq0$, both factors are units of $\overline K[t]$. If $\alpha\beta=0$, one factor is zero. In either case, $\bar{F}$ or $G$ is constant, contradicting independence.

Thus $Q_2$ has rank one and can be written over $K$ as $\kappa L(x,z)^2$, with $\kappa\in K^*$ and $L$ a nonzero linear form. After a simultaneous linear change of $x,z$, write
\[
 Q_2=\kappa x^2,\qquad \Lambda=\alpha x+\beta z.
\]
Then $\kappa f^2=\alpha f+\beta g$. Necessarily $\beta\neq0$, and~\eqref{eq:quadratic-form} becomes
\begin{equation}\label{eq:precanonical}
 H=\left(f(\rho),\ \beta z-2\kappa xf(\rho)+\alpha x,
 \ \frac{\kappa}{\beta}f(\rho)^2-\frac{\alpha}{\beta}f(\rho)\right),
 \qquad \rho=y+\kappa x^2.
\end{equation}
The shear $z'=z+(\alpha/\beta)x$ removes the terms involving $\alpha$. Finally set
\[
 x'=x,\qquad y'=y/\kappa,\qquad z''=(\beta/\kappa)z'.
\]
With $P(t)=f(\kappa t)$, the transformed map is $\Ncal_P(x',y',z'')$. All these changes are defined over $K$, and $P(0)=0$.
\end{proof}

\begin{proof}[Proof of Theorem~\ref{thm:classification}]
The forward implication follows from Theorem~\ref{thm:dependent-pair} and Proposition~\ref{prop:pair-normal}. For the converse, put $\rho=y+x^2$, $u=P(\rho)$, and $\bar{q}=P'(\rho)$. Then
\begin{equation}\label{eq:canonical-J}
 J\Ncal_P=
 \begin{pmatrix}
 2x\bar{q}&\bar{q}&0\\
 -2u-4x^2\bar{q}&-2x\bar{q}&1\\
 4xu\bar{q}&2u\bar{q}&0
 \end{pmatrix}.
\end{equation}
The trace and determinant vanish. The principal minors of order two are $2u\bar{q},0,-2u\bar{q}$, so nilpotence follows from~\eqref{eq:nilpotence}. The first and third components are independent because $u$ is nonconstant, and the second component has coefficient one in $z$. Hence all three are independent.
\end{proof}

\begin{remark}
Theorem~\ref{thm:classification} provides an affirmative answer to Problem~3.8 in~\cite{YanTang} and Problem~4.1 in \cite{HeYan}.
\end{remark}

\subsection{The dependent case over a polynomial coefficient ring}
To prove tameness, we also need a normal form for maps with dependent components. The coefficient ring is essential to the following elementary argument.

\begin{lemma}\label{lem:planar-form}
Let $R=K[z]$ and let $h=(h_1,h_2)\in R[x,y]^2$ have nilpotent planar Jacobian $J_{x,y}h$. There are $a_1,a_2,c_1,c_2\in R$ and $f\in R[t]$ such that
\begin{equation}\label{eq:dependent-form}
 h_1=a_2f(a_1x+a_2y)+c_1,\qquad
 h_2=-a_1f(a_1x+a_2y)+c_2,
\end{equation}
where $(a_1,a_2)=R$ and $f(0)=0$.
\end{lemma}

\begin{proof}
Set $c_i=h_i(0,0)$ and $\bar h=h-c$. Over $K(z)$, the normalized planar map $\bar h$ has nilpotent Jacobian and therefore dependent components. Clearing the denominators in a relation and dividing by the greatest common divisor gives
\[
 a_1\bar h_1+a_2\bar h_2=0,\qquad a_1,a_2\in R,\quad(a_1,a_2)=R.
\]
If $\bar h=0$, take $(a_1,a_2)=(1,0)$ and $f=0$. Otherwise, the primitive relation gives
$\bar h=(a_2q,-a_1q)$ for $q\in R[x,y]$. For example, if $u a_1+v a_2=1$, then $q=v\bar h_1-u\bar h_2$.

Choose $\alpha,\beta\in R$ with $\alpha a_1-\beta a_2=1$, and put
\begin{equation}\label{eq:unimodular-M}
 M(z)=\begin{pmatrix}a_1&a_2\\ \beta&\alpha\end{pmatrix}\in\SL_2(R),
 \qquad (s,t)=(a_1x+a_2y,\beta x+\alpha y).
\end{equation}
The trace equation is $a_2q_x-a_1q_y=0$. In the coordinates $s,t$, the derivation $a_2\partial_x-a_1\partial_y$ is $-\partial_t$. Hence $q=f(s)$ for $f\in R[s]$. Since $\bar h(0,0)=0$ and the pair is primitive, $f(0)=0$.
\end{proof}

\begin{remark}\label{rem:coefficient-ring}
Lemma~\ref{lem:planar-form} is Theorem 7.2.25 in \cite{EssenBook}. We give the proof here for completeness. In Lemma~\ref{lem:planar-form}, $f$ is a polynomial in one variable with coefficients in $K[z]$. The proof does not justify replacing $K[z][t]$ by $K[t]$. This distinction also applies to the dependent normal forms in Section~\ref{sec:blocks}.
\end{remark}

\subsection{Tame factorizations and polynomial inverses}
\begin{theorem}[Tameness]\label{thm1}
Let $H\in K[x,y,z]^3$ have nilpotent Jacobian. Then $F=\id+H$ is tame over $K$. More generally, for an indeterminate $s$,
\[
 \id+sH\in\TA_3(K[s]).
\]
\end{theorem}

\begin{proof}
Write $H=\bar H+c$, where $c=H(0)$. Then
\begin{equation}\label{eq:translation-normalize}
 \id+sH=\trans{sc}\circ(\id+s\bar H).
\end{equation}
It suffices to treat $\bar H$. We distinguish the cases of linearly independent and linearly dependent components after normalization.

Suppose first that the components of $\bar H$ are independent. By Theorem~\ref{thm:classification}, a constant conjugation transforms the map into $\Ncal_P$. Define
\[
 T_Q(x,y,z)=(x,y+Q(x,z),z),\qquad
 S_s(x,y,z)=(x+sf(y),y,z+sg(y)).
\]
For the quadratic-coordinate form~\eqref{eq:quadratic-form}, relation~\eqref{eq:quadratic-relation} gives
\begin{equation}\label{eq:tame-factorization}
 \id+sH=T_{Q-s\Lambda}^{-1}\circ S_s\circ T_Q.
\end{equation}
Indeed, the second component of the right-hand side is
\begin{align*}
 \rho-Q(x+sf(\rho),z+sg(\rho))+s\Lambda(x+sf(\rho),z+sg(\rho))
 &=y+s\bigl(-Q_xf(\rho)-Q_zg(\rho)+\Lambda(x,z)\bigr)\\
 &\quad+s^2\bigl(\Lambda(f(\rho),g(\rho))-Q_2(f(\rho),g(\rho))\bigr).
\end{align*}
The last line vanishes. The other two components follow directly from the composition. Each $T$ is elementary and $S_s$ is a product of two commuting elementary maps. In the canonical case one takes $Q=x^2$, $\Lambda=z$, $f=P$, and $g=P^2$.

Suppose next that the components of $\bar H$ are dependent. A constant conjugation gives $\bar H=(h_1,h_2,0)$. Apply Lemma~\ref{lem:planar-form} and let
\[
 S(x,y,z)=(M(z)(x,y)^{\mathsf T},z),
\]
where $M$ is given by~\eqref{eq:unimodular-M}. Since the third coordinate is fixed,
\begin{equation}\label{eq:dependent-tame-conjugate}
 S\circ(\id+s\bar H)\circ S^{-1}
 =\bigl(x+s(a_1c_1+a_2c_2),\
 y-sf(x)+s(\beta c_1+\alpha c_2),\ z\bigr).
\end{equation}
Here $f(x)\in K[z][x]$. This map is triangular and therefore tame over $K[s]$. The Euclidean algorithm over $K[z]$ expresses $M\in\SL_2(K[z])$ as a product of elementary matrices. Each factor lifts to an elementary automorphism of $K[s][x,y,z]$, so $S$ and $S^{-1}$ are tame. Undoing the constant conjugation and~\eqref{eq:translation-normalize} proves the theorem.
\end{proof}

The factorization~\eqref{eq:tame-factorization} also gives the inverse before the final linear normalization. For the quadratic-coordinate form~\eqref{eq:quadratic-form}, set
\begin{equation}\label{eq:quadratic-inverse}
 \begin{aligned}
 \pi&=Y+Q(X,Z)-s\Lambda(X,Z),\\
 x&=X-sf(\pi),\qquad z=Z-sg(\pi),\qquad y=\pi-Q(x,z).
 \end{aligned}
\end{equation}
Indeed, the inverse factorization is $T_Q^{-1}\circ S_{-s}\circ T_{Q-s\Lambda}$, which gives the displayed expressions. Thus~\eqref{eq:quadratic-inverse} is a two-sided polynomial inverse over $K[s]$.

In the canonical case, if
\[
 (X,Y,Z)=(x+sP(\rho),y+s(z-2xP(\rho)),z+sP(\rho)^2),\qquad \rho=y+x^2,
\]
put $\pi=Y+X^2-sZ$. Then
\begin{equation}\label{eq:canonical-inverse}
 (\id+s\Ncal_P)^{-1}(X,Y,Z)
 =\bigl(X-sP(\pi),\ \pi-(X-sP(\pi))^2,\ Z-sP(\pi)^2\bigr).
\end{equation}
Equation~\eqref{eq:intro-invariant} proves both compositions. The construction is polynomial in $s$ and involves no division by the parameter.

In treating dependent components, we also retain the constant term with respect to the planar variables. For example, $H=(z,x,0)$ vanishes at the origin and has nilpotent Jacobian, but $H((x,y,z)+sH(x,y,z))\neq H(x,y,z)$. Its constant term as a planar map over $K[z]$ is $(z,0)$. Direct substitution gives
\[
 (\id+sH)^{-1}(X,Y,Z)=(X-sZ,\ Y-sX+s^2Z,\ Z),
\]
which illustrates why normalization at the origin does not replace normalization in the two polynomial variables over the coefficient ring.

% \subsection{The earlier normal form and Problem 4.1}
% The normal form in Theorem~2.3 of the He--Yan manuscript~\cite{HeYan} requires a derivative in its second component. With this correction, the family takes the form
% \begin{equation}\label{eq:earlier-form}
%  \begin{aligned}
%  u&=g(ay+b(x)),\\
%  v&=\nu z-a^{-1}b'(x)u-\nu\ell x,\\
%  h&=c_0u^2+\ell u,\qquad
%  b(x)=\nu c_0ax^2+b_1x+b_0,
%  \end{aligned}
% \end{equation}
% where $a\nu c_0\neq0$. Normalization requires $g(b_0)=0$. If $b(x)$ is used in place of $b'(x)$, as in the printed formula, the Jacobian trace is
% \[
%  \bigl(b'(x)-b(x)\bigr)g'(ay+b(x)),
% \]
% which is generally nonzero. For~\eqref{eq:earlier-form}, the trace and determinant vanish and
% \[
%  \sigma_2(J(u,v,h))=u\,g'(ay+b(x))\bigl(b''(x)-2a\nu c_0\bigr)=0.
% \]
% The canonical family is the specialization $a=\nu=c_0=1$, $b=x^2$, $\ell=0$, and $g=P$. Thus Theorem~\ref{thm:classification} gives the proposed affirmative answer to Problem~4.1.

% The example following that problem does not contradict this conclusion. To verify this independently of the general classification, denote its nonzero parameters by $A,B,c$ and set
% \[
%  q=z-\frac{c}{2B}y^2,\qquad \delta=\frac{cB}{2A}.
% \]
% The printed map is $H=(Aq^2,Bq,cyq-\delta x)$. For the cyclic permutation $T(x,y,z)=(y,z,x)$,
% \begin{equation}\label{eq:example-cyclic}
%  T\circ H\circ T^{-1}
%  =\left(\rho,\frac cB x\rho-\delta z,\frac A{B^2}\rho^2\right),
%  \qquad \rho=By-\frac c2x^2.
% \end{equation}
% This is~\eqref{eq:earlier-form} with $a=B$, $\nu=-\delta$, $c_0=A/B^2$, $b(x)=-cx^2/2$, $g(t)=t$, and $\ell=0$. More explicitly, for
% \[
%  S(x,y,z)=\left(x,-\frac{2B}{c}y,\frac{B^2}{A}z\right),\qquad P(t)=-\frac c2t,
% \]
% we have $(S\circ T)\circ H\circ(S\circ T)^{-1}=\Ncal_P$. Its Jacobian is nilpotent and its components are independent; the asserted nonconjugacy is incorrect.


\subsection{The generic curve of the canonical family}\label{subsec:canonical-curve}
The canonical family gives an explicit model for the relative constant field used in Section~\ref{sec:dependent-pair}. Put
\[
 \rho=y+x^2,\qquad u=P(\rho),\qquad v=z-2xu.
\]
Since the third component of $\Ncal_P$ is $u^2$, the image field and its relative algebraic closure in $\mathcal K=K(x,y,z)$ are
\begin{equation}\label{eq:canonical-fields}
 E=K(u,v),\qquad L=K(\rho,v),\qquad \mathcal K=L(x),\qquad [L:E]=\deg P.
\end{equation}
Indeed, $y=\rho-x^2$ and $z=v+2xP(\rho)$, so $\mathcal K=L(x)$. The field $L$ is relatively algebraically closed in this purely transcendental extension. Moreover, $L/E$ is finite of degree $\deg P$, the degree of the univariate extension $K(\rho)/K(P(\rho))$ after adjoining the independent variable $v$.

Over $L$, the embedded generic curve has the parametrization
\begin{equation}\label{eq:canonical-conic}
 x\longmapsto\bigl(x,\ \rho-x^2,\ v+2P(\rho)x\bigr).
\end{equation}
It is a parabola in the containing affine plane, with the unique direction at infinity $[0:1:0]$. If $\deg P>1$, the geometric generic fiber over $E$ has several components, corresponding to the distinct generic roots of $P(\rho)=u$. Passing to $L$ selects one geometrically integral component. This explains the use of the relative algebraic closure in the proof of Proposition~\ref{prop:line-conic}.

The source coordinate $\rho$ and the target relation $W-U^2=0$ have different roles. The source coordinate describes the conic fibers, whereas the target relation defines the one-dimensional relation space used in the descent in Section~\ref{sec:dependent-pair}.
